\documentclass[11pt]{article}
\usepackage[margin=1in]{geometry}
\usepackage{amsmath,amssymbol,booktabs}
\title{End-to-end scores for the cascade: definition and testing}
\date{}
\begin{document}
\maketitle

\section*{Goal}
Given an observed cloud $x$ and the cascade's fit $(\hat f,\hat\theta)$, we ask whether the fitted model
reproduces $x$. We simulate patterns from the fit and compare them with $x$ using a proper scoring rule
(lower is better). All models are scored on the same cloud, so only the fit differs.

\section*{Scores}
\paragraph{Kernel score (primary).}
Let $\phi_i$ be the local configurations of $x$: the points within $R=1.5$ mean spacings of a centre (one of
$x$'s own points), translated to the origin and divided by $R$. Let $\Phi,\Phi'$ be configurations from
\emph{different} simulated patterns of model $P$. With a kernel $k$,
\[
  S(P,x) \;=\; \tfrac12\,\mathbb{E}\,k(\Phi,\Phi') \;-\; \frac1m\sum_{i=1}^{m}\mathbb{E}\,k(\phi_i,\Phi).
\]
For a characteristic $k$ this is strictly proper for the law of radius-$R$ configurations:
$\mathbb{E}_Q S(P,\cdot)=\tfrac12\|\mu_P-\mu_Q\|^2+\text{const}$. Each configuration is embedded as a
smoothed $21\times21$ grid ($h=0.1$), so $k$ is a Gaussian kernel of Euclidean distances between embeddings.
Its width is set once per cloud, to the median distance between $x$'s own configurations, and reused for
every model, which is what makes scores comparable. We use $m\le300$ centres and 16 simulations per model.

\paragraph{Dawid--Sebastiani score (secondary).}
With $T$ a vector of scale-free geometric statistics (Voronoi cell areas, Delaunay neighbour counts and angles;
the persistent-homology statistics are dropped to avoid overlap with the PersLay models being scored), and
$\mu_P,\Sigma_P$ the mean and shrunk covariance of $T$ over 150 simulations,
\[
  S(P,x)=\log\det\Sigma_P+(T(x)-\mu_P)^\top\Sigma_P^{-1}(T(x)-\mu_P).
\]
It is proper but not strictly (models agreeing on two moments tie).

\paragraph{Reference models and derived quantities.}
Two reference models are simulated per cloud: the \emph{oracle} (true family, true $\theta$) and \emph{CSR}
(Poisson with $\bar n=n$). Then
\[
  \text{regret}=S(\text{fit})-S(\text{oracle}),\qquad
  \text{gain}=S(\text{CSR})-S(\text{oracle}),\qquad
  \text{skill}=1-\frac{\sum\text{regret}}{\sum\text{gain}},
\]
pooled over clouds, and reported only where the mean gain exceeds 2 standard errors. Skill 1 means as good as
the truth; skill 0 means no better than CSR.

\section*{How the scores were tested}
\paragraph{Power check.}
On 300 test clouds (15 per family and $\tilde\delta$ bin, plus 60 Poisson clouds), only the oracle and CSR
were scored. A usable score should have gain $>0$ where structure is strong, and gain $\approx0$ on Poisson.
We report $d'=\overline{\text{gain}}/\text{sd}$ per cloud. The first score tried, a Wasserstein-1 energy score
between whole patterns, had no power; the two scores above are its replacements.

\begin{table}[h]
\centering
\begin{tabular}{lccccc}
\toprule
$d'$ & Thomas & nested & LGCP & Mat\'ern II & Poisson (floor) \\
\midrule
Kernel, $R{=}1.5$, point centres & 1.19 & 1.20 & 1.30 & 1.06 & 0.12 \\
DS, non-PH statistics            & 0.74 & 0.64 & 1.00 & 0.91 & 0.01 \\
Wasserstein-1 energy (failed)    & 0.64 & 0.00 & 0.86 & 0.04 & $-0.13$ \\
\bottomrule
\end{tabular}
\caption{Power check, $\tilde\delta\ge4$ bin (strongest structure); 15 clouds per family cell.}
\end{table}

\paragraph{Findings.}
(i) At $\tilde\delta\ge4$ the kernel score separates the truth from CSR for every family ($z=4.1$--$5.0$).
(ii) For $\tilde\delta<1$ it is mostly noise, and nested at $\tilde\delta<0$ goes the wrong way ($z=-2.1$);
DS is no better there. (iii) Other kernel settings were biased on Poisson clouds (e.g.\ $R=4$ with grid
centres, $z\approx-2.6$); the chosen setting was picked after seeing these results, on the same clouds.

\paragraph{End-to-end run.}
Scoring the assembled cascade on 300 clouds gave regret near zero and skill between 0.94 and 1.22. Skill above
1 means the fit scored better than the truth, so this is noise: regret standard errors are as large as the
regrets. The score cannot yet distinguish a good fit from the truth at this sample size.

\paragraph{Missing.}
There is no negative control, i.e.\ no test that a deliberately wrong fit (perturbed $\theta$, wrong family)
scores worse than the oracle. The power check only compares the truth with CSR, a large difference.

\end{document}
